\documentclass[11pt]{article}
\usepackage[margin=1in]{geometry}
\usepackage[UTF8,fontset=fandol]{ctex}
\usepackage{amsmath,amssymb,amsthm}
\usepackage{array}
\usepackage{parskip}
\usepackage{booktabs}

\newtheorem{theorem}{Theorem}
\newtheorem{lemma}[theorem]{Lemma}
\newtheorem{corollary}[theorem]{Corollary}
\newtheorem{proposition}[theorem]{Proposition}
\theoremstyle{definition}
\newtheorem{definition}[theorem]{Definition}
\theoremstyle{remark}
\newtheorem{remark}[theorem]{Remark}

\newcommand{\Z}{\mathbb{Z}}
\newcommand{\Sh}{\mathrm{Sh}}
\newcommand{\vtwo}{v_2}

\title{Disproof of conjecture \texttt{00000000437}}
\author{earthking11}
\date{2026-09-13}

\begin{document}
\maketitle

\section*{Verdict: FALSE}

We disprove conjecture \texttt{00000000437} as stated. The numerical formula in
the conjecture (Witt's formula) is a true classical theorem, but the appended
\emph{parity clause} is false. The refutation is unconditional and elementary:
two powers of two, namely $n = 1$ and $n = 2$, have different parities,
$L_1 = 2$ and $L_2 = 1$, so the parity of $L_n$ is not determined by whether
$n$ is a power of two.

\section{The conjecture}

Quoted verbatim from \texttt{conjectures/00000000437.md}:

\begin{quote}
\textbf{Definition:} The shuffle algebra $\Sh(V)$ (tensor algebra with shuffle
product). \textbf{Conjecture:} The $n$-th graded dimension of the
Lie-primitive space of $\Sh(V)$ for $\dim V = 2$ is
$\frac1n\sum_{d\mid n}\mu(d)\cdot 2^{n/d}$ (Witt's formula), and its parity is
determined explicitly by whether $n$ is a power of $2$.
\end{quote}

\begin{quote}
\textbf{中文。} 定义：shuffle 代数 $\Sh(V)$（张量代数带 shuffle 积）。猜想：
shuffle 代数 $\Sh(V)$ 的李代数本原空间的第 $n$ 个分次维数在 $\dim V = 2$
时为 $\frac1n\sum_{d\mid n}\mu(d)\cdot 2^{n/d}$（Witt 公式），其奇偶性由
$n$ 是否为 $2$ 的幂显式给出。
\end{quote}

For the rest of the document let
\begin{equation}\label{eq:witt}
  L_n \;=\; \frac1n \sum_{d\mid n} \mu(d)\, 2^{\,n/d},
\end{equation}
where $\mu$ is the Möbius function. Thus $L_n$ is the number of binary Lyndon
words of length $n$, equivalently the number of aperiodic binary necklaces of
length $n$, equivalently the number of monic irreducible polynomials of degree
$n$ over $\mathbb{F}_2$; all of these satisfy \eqref{eq:witt} (Witt's formula).
That formula is \emph{not} in dispute. The conjecture consists of two claims,
and only the second is false:
\begin{enumerate}
\item \textbf{(Witt's formula, true.)} The $n$-th graded dimension is $L_n$.
\item \textbf{(Parity clause, false.)} The parity of $L_n$ ``is determined
explicitly by whether $n$ is a power of $2$.''
\end{enumerate}

\section{The value table}

The values of $L_n$ and their parities for $n \le 20$ are as follows. The last
column records the predicate ``$n$ is a power of two'' (with the standard
convention $2^0 = 1$, so $n=1$ is a power of two).

\begin{center}
\begin{tabular}{r r l c}
\toprule
$n$ & $L_n$ & parity of $L_n$ & $n$ is a power of $2$ \\
\midrule
1  & 2       & even & yes \\
2  & 1       & odd  & yes \\
3  & 2       & even & no  \\
4  & 3       & odd  & yes \\
5  & 6       & even & no  \\
6  & 9       & odd  & no  \\
7  & 18      & even & no  \\
8  & 30      & even & yes \\
9  & 56      & even & no  \\
10 & 99      & odd  & no  \\
11 & 186     & even & no  \\
12 & 335     & odd  & no  \\
13 & 630     & even & no  \\
14 & 1161    & odd  & no  \\
15 & 2182    & even & no  \\
16 & 4080    & even & yes \\
17 & 7710    & even & no  \\
18 & 14532   & even & no  \\
19 & 27594   & even & no  \\
20 & 52377   & odd  & no  \\
\bottomrule
\end{tabular}
\end{center}

Reading the table, the parity of $L_n$ visibly fails to track the
power-of-two predicate: $n = 1, 2$ are both powers of two yet the parities are
even and odd respectively; $n = 8, 16$ are powers of two with even parity; and
$n = 6, 10, 12, 14$ are not powers of two yet have odd parity.

\section{The refutation}

\subsection{The sharpest witness: $n = 1$ and $n = 2$}

\begin{theorem}\label{thm:main}
$L_1 = 2$ and $L_2 = 1$. Both $1 = 2^0$ and $2 = 2^1$ are powers of two, but
$L_1$ is even and $L_2$ is odd. Hence there is no well-defined function of the
predicate ``$n$ is a power of two'' that returns the parity of $L_n$; the
parity clause of conjecture \texttt{00000000437} is false.
\end{theorem}

\begin{proof}
Directly from \eqref{eq:witt}. For $n = 1$ the only divisor is $d = 1$ and
$\mu(1) = 1$, so $L_1 = 1 \cdot 2^1 / 1 = 2$. For $n = 2$ the divisors are
$d = 1, 2$ with $\mu(1) = 1$, $\mu(2) = -1$, so
\[
  L_2 \;=\; \frac{\mu(1)2^2 + \mu(2)2^1}{2}
        \;=\; \frac{4 - 2}{2} \;=\; 1 .
\]
Thus $L_1 = 2 \equiv 0 \pmod 2$ while $L_2 = 1 \equiv 1 \pmod 2$. Since
$1 = 2^0$ and $2 = 2^1$ are both powers of two, the two arguments that the
parity clause says determine the same value in fact give different values.
\end{proof}

\begin{remark}[Under every faithful reading]
Read the clause as a statement about the truth value of the predicate $P(n)$ :=
``$n$ is a power of two''. Then $P(1) = P(2) = \text{true}$ while
$L_1 \bmod 2 \neq L_2 \bmod 2$, so no function of $P(n)$ can equal
$L_n \bmod 2$. The failure does not depend on which explicit formula the file
had in mind: the file supplies none, and none can exist, because the parity is
not a function of $P(n)$.
\end{remark}

\subsection{The fallback witness: $n = 2$ and $n = 8$}

If one declines to count $1$ as a power of two (some authors exclude $2^0$),
the refutation is unchanged, because the two \emph{proper} powers of two $2$ and
$8$ already have different parities:
\[
  L_2 = 1 \equiv 1 \pmod 2,
  \qquad
  L_8 = \frac{2^8 - 2^4}{8} = \frac{256 - 16}{8} = 30 \equiv 0 \pmod 2 .
\]
So even under the restricted convention ``powers of two are $2,4,8,16,\dots$'',
parity is not determined by whether $n$ is a power of two.

\subsection{The ``if and only if'' reading fails at $n = 6$}

A natural sharpening of the clause is the equivalence
\[
  \text{``$L_n$ is odd \quad iff \quad $n$ is a power of two''}.
\]
This is false in \emph{both} directions.
\begin{itemize}
\item $n = 8$ is a power of two, but $L_8 = 30$ is even. ($n = 16$ likewise:
$L_{16} = 4080$ is even.)
\item $n = 6$ is not a power of two, yet
\[
  L_6 = \frac{\mu(1)2^6 + \mu(2)2^3 + \mu(3)2^2 + \mu(6)2^1}{6}
      = \frac{64 - 8 - 4 + 2}{6} = \frac{54}{6} = 9
\]
is odd. The obstruction is intrinsic: $6 = 2^1 \cdot 3$ has $v_2(6) = 1$ and
odd part $3$, which is squarefree.
\end{itemize}
So neither implication holds, and the parity clause is false however the
unspecified ``explicit formula'' is read.

\section{The true repaired characterisation}

Although parity is \emph{not} governed by the power-of-two predicate alone, it
does have an exact description. It was verified numerically for all $n \le 60$
by \texttt{reproduce.py}; here is the proof.

\begin{proposition}\label{prop:repaired}
Write $n = 2^{a}\,m$ with $m$ odd and $a = \vtwo(n) \ge 0$. Then
\[
  L_n \text{ is odd}
  \quad\Longleftrightarrow\quad
  a \in \{1,2\} \ \text{ and }\ m \ \text{is squarefree}.
\]
Equivalently: $L_n$ is even for every odd $n$, and for even $n$ the value
$L_n$ is odd exactly when $\vtwo(n) \in \{1,2\}$ and the odd part of $n$ has no
repeated prime factor.
\end{proposition}

\begin{proof}
Multiplying \eqref{eq:witt} by $n$ gives the integer identity
\[
  n\,L_n \;=\; \sum_{d\mid n} \mu(d)\, 2^{\,n/d},
\]
whose nonzero summands are indexed by the squarefree divisors $d$ of $n$. On
that set the exponent $n/d$ is strictly decreasing in $d$, so there is a unique
squarefree divisor realising the smallest exponent, namely the largest one,
$d^\ast = \operatorname{rad}(n) = \prod_{p \mid n} p$. Every other summand is
divisible by $2^{\,n/d^\ast + 1}$, hence does not cancel the lowest-order term;
consequently
\[
  \vtwo\!\left(n L_n\right) \;=\; \frac{n}{\operatorname{rad}(n)} .
\]
Since $\vtwo(n) = a$ and $\vtwo(nL_n) = a + \vtwo(L_n)$, this gives
\[
  \vtwo(L_n) \;=\; \frac{n}{\operatorname{rad}(n)} - a .
\]
If $a = 0$ (that is, $n = m$ is odd) then $\operatorname{rad}(n) =
\operatorname{rad}(m)$ and
$\vtwo(L_n) = m/\operatorname{rad}(m) \ge 1$, so $L_n$ is even. If $a \ge 1$
then $\operatorname{rad}(n) = 2\operatorname{rad}(m)$ and
\[
  \vtwo(L_n) \;=\; 2^{\,a-1}\,\frac{m}{\operatorname{rad}(m)} - a .
\]
This vanishes exactly when (i) $m = \operatorname{rad}(m)$, i.e.\ $m$ is
squarefree, and (ii) $2^{\,a-1} = a$, i.e.\ $a \in \{1,2\}$: for $a \ge 3$ one
has $2^{\,a-1} > a$, so the expression is positive, and for $a \in \{1,2\}$ it
equals $m/\operatorname{rad}(m) - a$, which is $0$ precisely when $m$ is
squarefree. Hence $L_n$ is odd exactly for
$a \in \{1,2\}$ with $m$ squarefree.
\end{proof}

\begin{corollary}
The parity is governed by $\vtwo(n)$ \emph{and} the squarefreeness of the odd
part of $n$ \emph{together}. The power-of-two predicate alone is insufficient:
among the powers of two, $L_1, L_2, L_4$ are odd (for $n = 1,2,4$ with odd part
$1$ squarefree and $\vtwo(n) \in \{0,1,2\}$) while $L_8, L_{16}, L_{32}$ are
even ($\vtwo(n) \ge 3$). This is exactly the failure recorded in
Theorem~\ref{thm:main}.
\end{corollary}

\section{Caveats}

\begin{enumerate}
\item \textbf{Witt's formula is true and is not being disproved.} The
numerical part of the conjecture, $L_n = \frac1n\sum_{d\mid n}\mu(d)2^{n/d}$,
is a standard theorem (binary Lyndon words / aperiodic necklaces / irreducible
polynomials over $\mathbb{F}_2$). The verdict FALSE refers solely to the parity
clause.
\item \textbf{The parity clause is under-specified.} The file gives no explicit
formula for the parity, only the assertion that it ``is determined explicitly
by whether $n$ is a power of $2$.'' The refutation is robust: the sharpest
witness ($1,2$), the fallback ($2,8$), and the iff-violation ($n=6$) defeat,
respectively, the functional reading, the restricted-convention reading, and
the iff reading. No reading in which the parity is a function of the
power-of-two predicate alone can survive.
\item \textbf{Convention on $1$.} Under the convention $2^0 = 1$, the pair
$(1,2)$ is the smallest witness. Under the convention that $1$ is not a power
of two, the pair $(2,8)$ is the smallest witness. The verdict is the same.
\item \textbf{The repaired statement changes the hypothesis.} Replacing
``$n$ is a power of two'' by the pair of conditions in
Proposition~\ref{prop:repaired} yields a correct criterion; it is not a
cosmetic rewording of the conjecture, since $\vtwo(n) \in \{1,2\}$ and
squarefreeness of the odd part are independent of the power-of-two predicate.
For instance $n=6$ has odd $L_6$ yet is not a power of two, and $n=8$ is a
power of two with even $L_8$.
\item \textbf{Scope of the numerics.} The values $L_n$ for $n \le 60$ and the
repaired characterisation are checked by \texttt{reproduce.py}; the finite
witnesses $L_1, L_2, L_6, L_8, L_{24}$ and the parity statements are proved in
core Lean in \texttt{lean4/Main.lean}. The value table up to $n=20$ is quoted
in Section~3.
\end{enumerate}

\section{Reproducibility and formalisation}

The arithmetic is reproduced by \texttt{reproduce.py} (Python 3 standard
library only), which computes $\mu$ by factorisation, evaluates $L_n$ for
$n \le 60$, prints the table, asserts $L_1 = 2$, $L_2 = 1$, $L_6 = 9$, that
$1$ and $2$ are powers of two with different parities, the $(2,8)$ fallback,
the $n=6$ iff-violation, and the repaired characterisation; it prints
\texttt{PASS} and exits $0$ iff every check holds.

The finite witnesses are formalised in the Lean~4 project under
\texttt{lean4/} (core Lean only, \texttt{import Std}, no Mathlib, no
\texttt{sorry}). It proves, among others,
\begin{quote}\ttfamily
theorem witt\_1 : witt 1 = 2\\
theorem witt\_2 : witt 2 = 1\\
theorem witt\_6 : witt 6 = 9\\
theorem witt\_8 : witt 8 = 30\\
theorem isPow2\_1 : isPow2 1 = true\\
theorem isPow2\_2 : isPow2 2 = true\\
theorem isPow2\_6 : isPow2 6 = false\\
theorem witness\_1\_2\_parity : witt 1 \% 2 \(\neq\) witt 2 \% 2\\
theorem fallback\_2\_8\_parity : witt 2 \% 2 \(\neq\) witt 8 \% 2\\
theorem parity\_not\_determined\_by\_pow2 :\\
\hspace*{2em}\(\neg\) (\(\forall\) n m, isPow2 n = isPow2 m \(\to\) witt n \% 2 = witt m \% 2)\\
theorem conjecture\_00000000437\_false : \textit{(the collected refutation)}
\end{quote}
The Möbius function is computed from an \emph{ascending} smallest-prime-factor
search; this ordering is essential, since a ``largest divisor $\le \sqrt N$''
search returns the composite $4$ for $N = 20$ and $N = 24$ and silently
corrupts $\mu$ (for instance producing $L_{24} = 698869$ instead of the correct
$L_{24} = 698870$, proved as \texttt{witt\_24}). Every finite fact is closed by
\texttt{decide}, never \texttt{native\_decide}, so that no
\texttt{Lean.ofReduceBool} enters the axiom footprint; the
\texttt{\#print axioms} audit in \texttt{Check.lean} reports only
\texttt{propext} (and no axioms at all for the purely computational facts).

\section{Files}

\begin{center}
\begin{tabular}{l l}
\toprule
File & Description \\
\midrule
\texttt{README.md} & Verdict, bilingual quote, tables, witnesses, caveats, status. \\
\texttt{main.tex} & This document. Standalone \texttt{article}; \texttt{tectonic main.tex}. \\
\texttt{build/main.pdf} & PDF produced by \texttt{tectonic main.tex}. \\
\texttt{reproduce.py} & Standard-library reproduction; \texttt{PASS}/\texttt{FAIL}, non-zero on failure. \\
\texttt{lean4/lean-toolchain} & Pins \texttt{leanprover/lean4:v4.33.1}. \\
\texttt{lean4/lakefile.toml} & Lake configuration; library \texttt{Main}, package \texttt{tlmc437}. \\
\texttt{lean4/Main.lean} & Core-Lean formalisation of the witnesses. \\
\texttt{lean4/Check.lean} & \texttt{\#print axioms} audit. \\
\texttt{lean4/README.md} & Statement table and scope note. \\
\bottomrule
\end{tabular}
\end{center}

\end{document}
